\documentclass[11pt,fleqn, openany]{book} % Default font size and left-justified equations

%%%%%%%%%%%%%%%%%%%%%%%%%%%%%%%%%%%%%%%%%
% The Legrand Orange Book
% Structural Definitions File
% Version 2.1 (26/09/2018)
%
% Original author:
% Mathias Legrand (legrand.mathias@gmail.com) with modifications by:
% Vel (vel@latextemplates.com)
% 
% This file was downloaded from:
% http://www.LaTeXTemplates.com
%
% License:
% CC BY-NC-SA 3.0 (http://creativecommons.org/licenses/by-nc-sa/3.0/)
%
%%%%%%%%%%%%%%%%%%%%%%%%%%%%%%%%%%%%%%%%%

%----------------------------------------------------------------------------------------
%	VARIOUS REQUIRED PACKAGES AND CONFIGURATIONS
%----------------------------------------------------------------------------------------

\usepackage[table]{xcolor}

\usepackage{graphicx}
\usepackage{tabularx} % Required for including pictures
\usepackage{pgf,tikz,tkz-tab,eurosym,yhmath, stmaryrd}
\usepackage{pgfplots}
\usepackage{mathrsfs}
\usetikzlibrary{patterns}
\usetikzlibrary{trees}
\graphicspath{{../../Pictures/}}
\usepackage{multicol} 


\usepackage[english]{babel} % English language/hyphenation
\usepackage{icomma}
\usepackage{enumitem} % Customize lists
\setlist{nolistsep, nosep, nolistsep} % Reduce spacing between bullet points and numbered lists

\usepackage{booktabs} % Required for nicer horizontal rules in tables

 % Required for specifying colors by name


\definecolor{ocre}{RGB}{243,102,25} % Define the orange color used for highlighting throughout the book

\usepackage{listings}

\definecolor{codegreen}{rgb}{0,0.6,0}
\definecolor{codegray}{rgb}{0.5,0.5,0.5}
\definecolor{codepurple}{rgb}{0.58,0,0.82}
\definecolor{backcolour}{rgb}{0.95,0.95,0.92}

\lstdefinestyle{mystyle}{
    backgroundcolor=\color{backcolour},   
    commentstyle=\color{codegreen},
    keywordstyle=\color{magenta},
    numberstyle=\tiny\color{codegray},
    stringstyle=\color{codepurple},
    basicstyle=\ttfamily\footnotesize,
    breakatwhitespace=false,         
    breaklines=true,                 
    captionpos=b,                    
    keepspaces=true,                 
    numbers=left,                    
    numbersep=5pt,                  
    showspaces=false,                
    showstringspaces=false,
    showtabs=false,                  
    tabsize=2
}

\lstset{style=mystyle}

%----------------------------------------------------------------------------------------
% Paramétrage XSIM
%----------------------------------------------------------------------------------------

\usepackage[no-files]{xsim}


\DeclareExerciseEnvironmentTemplate{myex}{%
    \textbf{%
      \hypertarget{ex:\ExerciseID}{\sffamily{\ensuremath{\blacktriangleright}} Exercice \GetExerciseProperty{counter} \GetExerciseProperty{subtitle} --}
      \hyperlink{sol:\ExerciseID}{Voir le corrigé}%
    }\par
}{\par\smallskip}

\DeclareExerciseEnvironmentTemplate{mysol}{%
    \textbf{%
      \hypertarget{sol:\ExerciseID}{\sffamily{\ensuremath{\blacktriangleright}} Correction \GetExerciseProperty{counter} --}
      \hyperlink{ex:\ExerciseID}{Voir l'énoncé}%
    }\par
}{\par\medskip}

\xsimsetup{
  exercise/template = myex ,
  solution/template = mysol 
}

%Collection exercices

\DeclareExerciseTagging{topic}

\xsimsetup{collect}

%----------------------------------------------------------------------------------------
% SYMBOLES
%----------------------------------------------------------------------------------------

\newcommand\imCMsym[4][\mathord]{%
  \DeclareFontFamily{U} {#2}{}
  \DeclareFontShape{U}{#2}{m}{n}{
    <-6> #25
    <6-7> #26
    <7-8> #27
    <8-9> #28
    <9-10> #29
    <10-12> #210
    <12-> #212}{}
  \DeclareSymbolFont{CM#2} {U} {#2}{m}{n}
  \DeclareMathSymbol{#4}{#1}{CM#2}{#3}
}
\newcommand\alsoimCMsym[4][\mathord]{\DeclareMathSymbol{#4}{#1}{CM#2}{#3}}

\imCMsym{cmmi}{124}{\CMjmath}

\newcommand{\Oij}{(O\,;\,\vec{\imath}\,,\, \vec{\CMjmath} )}
\newcommand{\Oijk}{(O\,;\,\vec{\imath}\,,\, \vec{\CMjmath}\,,\,\vec{k})}

\newcommand\e{\mathrm{e}}
\newcommand\R{\mathbb{R}}
\newcommand\N{\mathbb{N}}


%----------------------------------------------------------------------------------------
%	MARGINS
%----------------------------------------------------------------------------------------

\usepackage{geometry} % Required for adjusting page dimensions and margins

\geometry{
	paper=a4paper, % Paper size, change to letterpaper for US letter size
	top=3cm, % Top margin
	bottom=3cm, % Bottom margin
	left=2cm, % Left margin
	right=2cm, % Right margin
	headheight=14pt, % Header height
	footskip=1.4cm, % Space from the bottom margin to the baseline of the footer
	headsep=10pt, % Space from the top margin to the baseline of the header
	%showframe, % Uncomment to show how the type block is set on the page
}

\setlength{\parindent}{0pt}
\parskip=5pt



%----------------------------------------------------------------------------------------
%	FONTS
%----------------------------------------------------------------------------------------

\usepackage{avant} % Use the Avantgarde font for headings
\usepackage{times} % Use the Times font for headings
\usepackage{mathptmx} % Use the Adobe Times Roman as the default text font together with math symbols from the Sym­bol, Chancery and Com­puter Modern fonts

%\usepackage{microtype} % Slightly tweak font spacing for aesthetics
%\usepackage[utf8]{inputenc} % Required for including letters with accents
\usepackage[T1]{fontenc} % Use 8-bit encoding that has 256 glyphs

%----------------------------------------------------------------------------------------
%	BIBLIOGRAPHY AND INDEX
%----------------------------------------------------------------------------------------

\usepackage[style=numeric,citestyle=numeric,sorting=nyt,sortcites=true,autopunct=true,babel=hyphen,hyperref=true,abbreviate=false,backref=true,backend=biber]{biblatex}
\addbibresource{bibliography.bib} % BibTeX bibliography file
\defbibheading{bibempty}{}

\usepackage{calc} % For simpler calculation - used for spacing the index letter headings correctly
\usepackage{makeidx} % Required to make an index
\makeindex % Tells LaTeX to create the files required for indexing

%----------------------------------------------------------------------------------------
%	MAIN TABLE OF CONTENTS
%----------------------------------------------------------------------------------------

\usepackage{titletoc} % Required for manipulating the table of contents

\contentsmargin{0cm} % Removes the default margin

% Part text styling (this is mostly taken care of in the PART HEADINGS section of this file)
\titlecontents{part}
	[0cm] % Left indentation
	{\addvspace{20pt}\bfseries} % Spacing and font options for parts
	{}
	{}
	{}

% Chapter text styling
\titlecontents{chapter}
	[1.25cm] % Left indentation
	{\addvspace{12pt}\large\sffamily\bfseries} % Spacing and font options for chapters
	{\color{ocre!60}\contentslabel[\Large\thecontentslabel]{1.25cm}\color{ocre}} % Formatting of numbered sections of this type
	{\color{ocre}} % Formatting of numberless sections of this type
	{\color{ocre!60}\normalsize\;\titlerule*[.5pc]{.}\;\thecontentspage} % Formatting of the filler to the right of the heading and the page number

% Section text styling
\titlecontents{section}
	[1.25cm] % Left indentation
	{\addvspace{3pt}\sffamily\bfseries} % Spacing and font options for sections
	{\contentslabel[\thecontentslabel]{1.25cm}} % Formatting of numbered sections of this type
	{} % Formatting of numberless sections of this type
	{\hfill\color{black}\thecontentspage} % Formatting of the filler to the right of the heading and the page number

% Subsection text styling
\titlecontents{subsection}
	[1.25cm] % Left indentation
	{\addvspace{1pt}\sffamily\small} % Spacing and font options for subsections
	{\contentslabel[\thecontentslabel]{1.25cm}} % Formatting of numbered sections of this type
	{} % Formatting of numberless sections of this type
	{\ \titlerule*[.5pc]{.}\;\thecontentspage} % Formatting of the filler to the right of the heading and the page number

% Figure text styling
\titlecontents{figure}
	[1.25cm] % Left indentation
	{\addvspace{1pt}\sffamily\small} % Spacing and font options for figures
	{\thecontentslabel\hspace*{1em}} % Formatting of numbered sections of this type
	{} % Formatting of numberless sections of this type
	{\ \titlerule*[.5pc]{.}\;\thecontentspage} % Formatting of the filler to the right of the heading and the page number

% Table text styling
\titlecontents{table}
	[1.25cm] % Left indentation
	{\addvspace{1pt}\sffamily\small} % Spacing and font options for tables
	{\thecontentslabel\hspace*{1em}} % Formatting of numbered sections of this type
	{} % Formatting of numberless sections of this type
	{\ \titlerule*[.5pc]{.}\;\thecontentspage} % Formatting of the filler to the right of the heading and the page number

%----------------------------------------------------------------------------------------
%	MINI TABLE OF CONTENTS IN PART HEADS
%----------------------------------------------------------------------------------------

% Chapter text styling
\titlecontents{lchapter}
	[0em] % Left indentation
	{\addvspace{15pt}\large\sffamily\bfseries} % Spacing and font options for chapters
	{\color{ocre}\contentslabel[\Large\thecontentslabel]{1.25cm}\color{ocre}} % Chapter number
	{}  
	{\color{ocre}\normalsize\sffamily\bfseries\;\titlerule*[.5pc]{.}\;\thecontentspage} % Page number

% Section text styling
\titlecontents{lsection}
	[0em] % Left indentation
	{\sffamily\small} % Spacing and font options for sections
	{\contentslabel[\thecontentslabel]{1.25cm}} % Section number
	{}
	{}

% Subsection text styling (note these aren't shown by default, display them by searchings this file for tocdepth and reading the commented text)
\titlecontents{lsubsection}
	[.5em] % Left indentation
	{\sffamily\footnotesize} % Spacing and font options for subsections
	{\contentslabel[\thecontentslabel]{1.25cm}}
	{}
	{}

%----------------------------------------------------------------------------------------
%	HEADERS AND FOOTERS
%----------------------------------------------------------------------------------------


\usepackage{fancyhdr} % Required for header and footer configuration

\pagestyle{fancy}
\renewcommand{\chaptermark}[1]{\markboth{\sffamily\normalsize\bfseries\ \thechapter.\ #1}{}} % Chapter text font settings
\renewcommand{\sectionmark}[1]{\markright{\sffamily\normalsize\thesection\hspace{5pt}#1}{}} % Section text font settings
\fancyhf{} \fancyhead[LE,RO]{\sffamily\normalsize\thepage} % Font setting for the page number in the header
\fancyhead[LO]{\rightmark} % Print the nearest section name on the left side of odd pages
\fancyhead[RE]{\leftmark} % Print the current chapter name on the right side of even pages

\fancyfoot[L]{Jason LAPEYRONNIE}
\fancyfoot[R]{\href{http://mathoutils.fr}{http://mathoutils.fr}} % Uncomment to include a footer

\renewcommand{\headrulewidth}{0.5pt} % Thickness of the rule under the header
\renewcommand{\footrulewidth}{0.5pt} % Thickness of the rule under the header

\fancypagestyle{plain}{% Style for when a plain pagestyle is specified
	\fancyhead{}\renewcommand{\headrulewidth}{0pt}%
}

% Removes the header from odd empty pages at the end of chapters
\makeatletter
\renewcommand{\cleardoublepage}{
\clearpage\ifodd\c@page\else
\hbox{}
\vspace*{\fill}
\thispagestyle{empty}
\newpage
\fi}

%----------------------------------------------------------------------------------------
%	THEOREM STYLES
%----------------------------------------------------------------------------------------

\usepackage{amsmath,amsfonts,amssymb,amsthm} % For math equations, theorems, symbols, etc

\newcommand{\intoo}[2]{\mathopen{]}#1\,;#2\mathclose{[}}
\newcommand{\ud}{\mathop{\mathrm{{}d}}\mathopen{}}
\newcommand{\intff}[2]{\mathopen{[}#1\,;#2\mathclose{]}}
\renewcommand{\qedsymbol}{$\blacksquare$}
\newtheorem{notation}{Notation}[section]

% Boxed/framed environments
\newtheoremstyle{ocrenumbox}% Theorem style name
{0pt}% Space above
{0pt}% Space below
{\normalfont}% Body font
{}% Indent amount
{\small\bf\sffamily\color{ocre}}% Theorem head font
{\;:\;}% Punctuation after theorem head
{0.25em}% Space after theorem head
{\small\sffamily\color{ocre}\thmname{#1}\nobreakspace\thmnumber{\@ifnotempty{#1}{}\@upn{#2}}% Theorem text (e.g. Theorem 2.1)
\thmnote{\nobreakspace\the\thm@notefont\sffamily\bfseries\color{black}---\nobreakspace#3}} % Optional theorem note

\newtheoremstyle{blacknumex}% Theorem style name
{5pt}% Space above
{10pt}% Space below
{\normalfont}% Body font
{} % Indent amount
{\small\bf\sffamily}% Theorem head font
{\;:\;}% Punctuation after theorem head
{0.25em}% Space after theorem head
{\small\sffamily{\tiny\ensuremath{\blacksquare}}\nobreakspace\thmname{#1}\nobreakspace\thmnumber{\@ifnotempty{#1}{}\@upn{#2}}% Theorem text (e.g. Theorem 2.1)
\thmnote{\nobreakspace\the\thm@notefont\sffamily\bfseries---\nobreakspace#3}}% Optional theorem note

\newtheoremstyle{blacknumexo}% Theorem style name
{15pt}% Space above
{10pt}% Space below
{\normalfont}% Body font
{} % Indent amount
{\small\bf\sffamily}% Theorem head font
{}% Punctuation after theorem head
{0.5em}% Space after theorem head
{\small\sffamily{\ensuremath{\blacktriangleright}}\nobreakspace\thmname{#1}\nobreakspace\thmnumber{\@ifnotempty{#1}{}\@upn{#2}}% Theorem text (e.g. Theorem 2.1)
\thmnote{\nobreakspace\the\thm@notefont\sffamily\bfseries---\nobreakspace#3} \\}% Optional theorem note



\newtheoremstyle{blacknumbox} % Theorem style name
{0pt}% Space above
{5pt}% Space below
{}% Body font
{}% Indent amount
{\large\bf\sffamily}% Theorem head font
{\;:\;}% Punctuation after theorem head
{0.25em}% Space after theorem head
{\small\sffamily\thmname{#1}\nobreakspace\thmnumber{\@ifnotempty{#1}{}\@upn{#2}}% Theorem text (e.g. Theorem 2.1)
\thmnote{\nobreakspace\the\thm@notefont\sffamily\bfseries---\nobreakspace#3}}% Optional theorem note

% Non-boxed/non-framed environments
\newtheoremstyle{ocrenum}% Theorem style name
{5pt}% Space above
{5pt}% Space below
{\normalfont}% Body font
{}% Indent amount
{\small\bf\sffamily\color{ocre}}% Theorem head font
{\;:\;}% Punctuation after theorem head
{0.25em}% Space after theorem head
{\small\sffamily\color{ocre}\thmname{#1}\nobreakspace\thmnumber{\@ifnotempty{#1}{}\@upn{#2}}% Theorem text (e.g. Theorem 2.1)
\thmnote{\nobreakspace\the\thm@notefont\sffamily\bfseries\color{black}---\nobreakspace#3}} % Optional theorem note
\makeatother

% Defines the theorem text style for each type of theorem to one of the three styles above
\newcounter{dummy} 
\newcounter{thm}
\newcounter{correction}
\newcounter{qst}
\theoremstyle{ocrenumbox}
\newtheorem{theoremeT}[dummy]{Théorème}
\newtheorem{exerciseT}{Propriété}
\newtheorem{principeT}{Principe}
\theoremstyle{blacknumex}
\newtheorem{exampleT}{Exemple}
\theoremstyle{blacknumexo}
\newtheorem{exo}[thm]{Exercice}
\newtheorem{corr}[correction]{Correction}
\newtheorem{quest}[qst]{Question}
\theoremstyle{blacknumbox}
\newtheorem{vocabulary}{Vocabulary}[section]
\newtheorem{definitionT}{Définition}
\newtheorem{corollaryT}[dummy]{Corollary}
\theoremstyle{ocrenum}
\newtheorem{proofT}[dummy]{Démonstration}


%----------------------------------------------------------------------------------------
%	DEFINITION OF COLORED BOXES
%----------------------------------------------------------------------------------------

\RequirePackage[framemethod=default]{mdframed} % Required for creating the theorem, definition, exercise and corollary boxes

% Theorem box
\newmdenv[skipabove=7pt,
skipbelow=7pt,
backgroundcolor=black!5,
linecolor=ocre,
innerleftmargin=5pt,
innerrightmargin=5pt,
innertopmargin=10pt,
leftmargin=0cm,
rightmargin=0cm,
innerbottommargin=5pt]{tBox}

%Proposition box	  
\newmdenv[skipabove=7pt,
skipbelow=7pt,
rightline=false,
leftline=true,
topline=false,
bottomline=false,
backgroundcolor=ocre!10,
linecolor=ocre,
innerleftmargin=5pt,
innerrightmargin=5pt,
innertopmargin=10pt,
innerbottommargin=3pt,
leftmargin=0cm,
rightmargin=0cm,
linewidth=4pt]{eBox}	

% Definition box
\newmdenv[skipabove=7pt,
backgroundcolor=ocre!4,
skipbelow=7pt,
rightline=false,
leftline=true,
topline=false,
bottomline=false,
linecolor=ocre,
innerleftmargin=5pt,
innerrightmargin=5pt,
innertopmargin=10pt,
leftmargin=0cm,
rightmargin=0cm,
linewidth=4pt,
innerbottommargin=5pt]{dBox}	

% Corollary box
\newmdenv[skipabove=7pt,
skipbelow=7pt,
rightline=false,
leftline=true,
topline=false,
bottomline=false,
linecolor=gray,
backgroundcolor=black!5,
innerleftmargin=5pt,
innerrightmargin=5pt,
innertopmargin=5pt,
leftmargin=0cm,
rightmargin=0cm,
linewidth=4pt,
innerbottommargin=5pt]{cBox}

\newmdenv[skipabove=7pt,
skipbelow=7pt,
backgroundcolor=black!5,
innerleftmargin=5pt,
topline=false,
bottomline=false,
rightline=false,
leftline=false,
innerrightmargin=5pt,
innertopmargin=5pt,
leftmargin=0cm,
rightmargin=0cm,
innerbottommargin=5pt]{xBox}

% Creates an environment for each type of theorem and assigns it a theorem text style from the "Theorem Styles" section above and a colored box from above
\newenvironment{theorem}{\begin{tBox}\begin{theoremeT}}{\end{theoremeT}\end{tBox}}

\newenvironment{exo2}{\noindent \begin{exo}\item\relax \noindent \begin{eBox}\item\relax}{\end{eBox}\end{exo}}


\newenvironment{proposition}{\begin{eBox}\begin{exerciseT}}{\hfill{\color{ocre}}\end{exerciseT}\end{eBox}}		

\newenvironment{principe}{\begin{eBox}\begin{principeT}}{\hfill{\color{ocre}}\end{principeT}\end{eBox}}	
		  
\newenvironment{definition}{\begin{dBox}\begin{definitionT}}{\end{definitionT}\end{dBox}}	

\newenvironment{example}{\begin{xBox}\begin{exampleT}}{\hfill{\tiny\ensuremath{\blacksquare}}\end{exampleT}\end{xBox}}

\newenvironment{demonstration}{\begin{proofT}}{\hfill{\tiny\ensuremath{\square}}\end{proofT}}		
\newenvironment{corollary}{\begin{cBox}\begin{corollaryT}}{\end{corollaryT}\end{cBox}}	

%----------------------------------------------------------------------------------------
%	REMARK ENVIRONMENT
%----------------------------------------------------------------------------------------

\newenvironment{remark}{\par\vspace{5pt}\small % Vertical white space above the remark and smaller font size
\begin{list}{}{
\leftmargin=25pt % Indentation on the left
\rightmargin=15pt}\item\ignorespaces % Indentation on the right
\makebox[-2.5pt]{\begin{tikzpicture}[overlay]
\node[draw=ocre!60,line width=1pt,circle,fill=ocre!25,font=\sffamily\bfseries,inner sep=2pt,outer sep=0pt] at (-15pt,0pt){\textcolor{ocre}{R}};\end{tikzpicture}} % Orange R in a circle
\advance\baselineskip -1pt}{\end{list}\vskip5pt} % Tighter line spacing and white space after remark

%----------------------------------------------------------------------------------------
%	SECTION NUMBERING IN THE MARGIN
%----------------------------------------------------------------------------------------

\makeatletter
\renewcommand{\@seccntformat}[1]{\llap{\textcolor{ocre}{\csname the#1\endcsname}\hspace{1em}}}                    
\renewcommand{\section}{\@startsection{section}{1}{\z@}
{-4ex \@plus -1ex \@minus -.4ex}
{1ex \@plus.2ex }
{\normalfont\LARGE\sffamily\bfseries}}
\renewcommand{\subsection}{\@startsection {subsection}{2}{\z@}
{-3ex \@plus -0.1ex \@minus -.4ex}
{0.5ex \@plus.2ex }
{\normalfont\sffamily\bfseries}}
\renewcommand{\subsubsection}{\@startsection {subsubsection}{3}{\z@}
{-2ex \@plus -0.1ex \@minus -.2ex}
{.2ex \@plus.2ex }
{\normalfont\small\sffamily\bfseries}}                        
\renewcommand\paragraph{\@startsection{paragraph}{4}{\z@}
{-2ex \@plus-.2ex \@minus .2ex}
{.1ex}
{\normalfont\small\sffamily\bfseries}}

%----------------------------------------------------------------------------------------
%	PART HEADINGS
%----------------------------------------------------------------------------------------

% Numbered part in the table of contents
\newcommand{\@mypartnumtocformat}[2]{%
	\setlength\fboxsep{0pt}%
	\noindent\colorbox{ocre!20}{\strut\parbox[c][.7cm]{\ecart}{\color{ocre!70}\Large\sffamily\bfseries\centering#1}}\hskip\esp\colorbox{ocre!40}{\strut\parbox[c][.7cm]{\linewidth-\ecart-\esp}{\Large\sffamily\centering#2}}%
}

% Unnumbered part in the table of contents
\newcommand{\@myparttocformat}[1]{%
	\setlength\fboxsep{0pt}%
	\noindent\colorbox{ocre!40}{\strut\parbox[c][.7cm]{\linewidth}{\Large\sffamily\centering#1}}%
}

\newlength\esp
\setlength\esp{4pt}
\newlength\ecart
\setlength\ecart{1.2cm-\esp}
\newcommand{\thepartimage}{}%
\newcommand{\partimage}[1]{\renewcommand{\thepartimage}{#1}}%
\def\@part[#1]#2{%
\ifnum \c@secnumdepth >-2\relax%
\refstepcounter{part}%
\addcontentsline{toc}{part}{\texorpdfstring{\protect\@mypartnumtocformat{\thepart}{#1}}{\partname~\thepart\ ---\ #1}}
\else%
\addcontentsline{toc}{part}{\texorpdfstring{\protect\@myparttocformat{#1}}{#1}}%
\fi%
\startcontents%
\markboth{}{}%
{\thispagestyle{empty}%
\begin{tikzpicture}[remember picture,overlay]%
\node at (current page.north west){\begin{tikzpicture}[remember picture,overlay]%	
\fill[ocre!20](0cm,0cm) rectangle (\paperwidth,-\paperheight);
\node[anchor=north] at (4cm,-3.25cm){\color{ocre!40}\fontsize{220}{100}\sffamily\bfseries\thepart}; 
\node[anchor=south east] at (\paperwidth-1cm,-\paperheight+1cm){\parbox[t][][t]{8.5cm}{
\printcontents{l}{0}{\setcounter{tocdepth}{1}}% The depth to which the Part mini table of contents displays headings; 0 for chapters only, 1 for chapters and sections and 2 for chapters, sections and subsections
}};
\node[anchor=north east] at (\paperwidth-1.5cm,-3.25cm){\parbox[t][][t]{15cm}{\strut\raggedleft\color{white}\fontsize{30}{30}\sffamily\bfseries#2}};
\end{tikzpicture}};
\end{tikzpicture}}%
\@endpart}
\def\@spart#1{%
\startcontents%
\phantomsection
{\thispagestyle{empty}%
\begin{tikzpicture}[remember picture,overlay]%
\node at (current page.north west){\begin{tikzpicture}[remember picture,overlay]%	
\fill[ocre!20](0cm,0cm) rectangle (\paperwidth,-\paperheight);
\node[anchor=north east] at (\paperwidth-1.5cm,-3.25cm){\parbox[t][][t]{15cm}{\strut\raggedleft\color{white}\fontsize{30}{30}\sffamily\bfseries#1}};
\end{tikzpicture}};
\end{tikzpicture}}
\addcontentsline{toc}{part}{\texorpdfstring{%
\setlength\fboxsep{0pt}%
\noindent\protect\colorbox{ocre!40}{\strut\protect\parbox[c][.7cm]{\linewidth}{\Large\sffamily\protect\centering #1\quad\mbox{}}}}{#1}}%
\@endpart}
\def\@endpart{\vfil\newpage
\if@twoside
\if@openright
\null
\thispagestyle{empty}%
\newpage
\fi
\fi
\if@tempswa
\twocolumn
\fi}

%----------------------------------------------------------------------------------------
%	CHAPTER HEADINGS
%----------------------------------------------------------------------------------------

% A switch to conditionally include a picture, implemented by Christian Hupfer
\newif\ifusechapterimage
\usechapterimagetrue
\newcommand{\thechapterimage}{}%
\newcommand{\chapterimage}[1]{\ifusechapterimage\renewcommand{\thechapterimage}{#1}\fi}%
\newcommand{\autodot}{.}
\def\@makechapterhead#1{%
{\parindent \z@ \raggedright \normalfont
\ifnum \c@secnumdepth >\m@ne
\if@mainmatter
\begin{tikzpicture}[remember picture,overlay]
\node at (current page.north west)
{\begin{tikzpicture}[remember picture,overlay]
\node[anchor=north west,inner sep=0pt] at (0,0) {\ifusechapterimage\includegraphics[width=\paperwidth]{\thechapterimage}\fi};
\draw[anchor=west] (\Gm@lmargin,-3cm) node [line width=2pt,rounded corners=15pt,draw=ocre,fill=white,fill opacity=0.5,inner sep=15pt]{\strut\makebox[22cm]{}};
\draw[anchor=west] (\Gm@lmargin+.3cm,-3cm) node {\huge\sffamily\bfseries\color{black}\thechapter\autodot~#1\strut};
\end{tikzpicture}};
\end{tikzpicture}
\else
\begin{tikzpicture}[remember picture,overlay]
\node at (current page.north west)
{\begin{tikzpicture}[remember picture,overlay]
\node[anchor=north west,inner sep=0pt] at (0,0) {\ifusechapterimage\includegraphics[width=\paperwidth]{\thechapterimage}\fi};
\draw[anchor=west] (\Gm@lmargin,-3cm) node [line width=2pt,rounded corners=15pt,draw=ocre,fill=white,fill opacity=0.5,inner sep=15pt]{\strut\makebox[22cm]{}};
\draw[anchor=west] (\Gm@lmargin+.3cm,-3cm) node {\huge\sffamily\bfseries\color{black}#1\strut};
\end{tikzpicture}};
\end{tikzpicture}
\fi\fi\par\vspace*{50\p@}}}

%-------------------------------------------

\def\@makeschapterhead#1{%
\begin{tikzpicture}[remember picture,overlay]
\node at (current page.north west)
{\begin{tikzpicture}[remember picture,overlay]
\node[anchor=north west,inner sep=0pt] at (0,0) {\ifusechapterimage\includegraphics[width=\paperwidth]{\thechapterimage}\fi};
\draw[anchor=west] (\Gm@lmargin,-3cm) node [line width=2pt,rounded corners=15pt,draw=ocre,fill=white,fill opacity=0.5,inner sep=15pt]{\strut\makebox[22cm]{}};
\draw[anchor=west] (\Gm@lmargin+.3cm,-3cm) node {\huge\sffamily\bfseries\color{black}#1\strut};
\end{tikzpicture}};
\end{tikzpicture}
\par\vspace*{50\p@}}
\makeatother

%----------------------------------------------------------------------------------------
%	LINKS
%----------------------------------------------------------------------------------------

\usepackage{hyperref}
\hypersetup{hidelinks,backref=true,pagebackref=true,hyperindex=true,colorlinks=false,breaklinks=true,urlcolor=ocre,bookmarks=true,bookmarksopen=false}

\usepackage{bookmark}
\bookmarksetup{
open,
numbered,
addtohook={%
\ifnum\bookmarkget{level}=0 % chapter
\bookmarksetup{bold}%
\fi
\ifnum\bookmarkget{level}=-1 % part
\bookmarksetup{color=ocre,bold}%
\fi
}
}

\renewcommand*\thesection{\arabic{section}}

\newcommand*{\coord}[3]{% 
  \ensuremath{\overrightarrow{#1}\, 
    \begin{pmatrix} 
      #2\\ 
      #3 
    \end{pmatrix}}}
    
  \newcommand*{\coordb}[2]{% 
  \ensuremath{ 
    \begin{pmatrix} 
      #1\\ 
      #2 
    \end{pmatrix}}}

\newcommand*{\coorde}[4]{% 
  \renewcommand{\arraystretch}{1}\ensuremath{\overrightarrow{#1}\, 
    \begin{pmatrix} 
      #2\\ 
      #3 \\
      #4
    \end{pmatrix}}}    
  \newcommand*{\coordbe}[3]{% 
 \renewcommand{\arraystretch}{1} \ensuremath{ 
    \begin{pmatrix} 
      #1\\ 
      #2 \\
      #3
    \end{pmatrix}}}  
    
\newcommand{\Card}{\mathrm{Card}}




\begin{document}

\chapterimage{../../Pictures/background}


\chapter{Cours : Primitives et équations différentielles}


\section{Notion d'équation différentielle}

\begin{definition}[Équation différentielle]Une équation différentielle est une équation dont l'inconnue est une fonction notée $y$ et qui fait intervenir les dérivées de cette fonction.\end{definition}

\begin{example}On considère la fonction $f$ définie pour tout réel $x$ par $f(x)=\e^{4x-2}$. 

$f$ est solution de l'équation différentielle $y'=4y$. 

En effet, $f$ est dérivable sur $\mathbb{R}$ et pour tout réel $x$, $f'(x)=4\e^{4x-2}$, c'est-à-dire $f'(x)=4f(x)$.\end{example}

\begin{example}On considère la fonction $f$ définie pour tout réel $x$ par $f(x)=\ln(1+\e^x)\e^{-x}$. 

$f$ est solution de l'équation différentielle $y'+y=\dfrac{1}{1+\e^x}$. En effet :

\begin{itemize}
\item Pour tout réel $x$, on pose $u(x)=\ln(1+\e^x)$. $u$ est définie et dérivable sur $\mathbb{R}$ et pour tout réel $x$, \[u'(x)=\dfrac{\e^x}{1+\e^x}.\]
\item Pour tout réel $x$, on pose $v(x)=\e^{-x}$. $v$ est dérivable sur $\mathbb{R}$ et pour tout réel $x$, $v'(x)=-\e^{-x}$.
\end{itemize}
Ainsi, $f$ est dérivable sur $\mathbb{R}$ et pour tout réel $x$,
\[f'(x)=\dfrac{\e^x}{1+\e^x}\times \e^{-x}+\ln(1+\e^x)\times(-\e^{-x})=\dfrac{1}{1+\e^x}-\ln(1+\e^x)\e^{-x}.\]
Ainsi, pour tout réel $x$,
\[f'(x)+f(x)=\dfrac{1}{1+\e^x}-\ln(1+\e^x)\e^{-x}+\ln(1+\e^x)\e^{-x}=\dfrac{1}{1+\e^x}.\]
$f$ est donc bien solution de l'équation différentielle $y'+y=\dfrac{1}{1+\e^x}$.\end{example}


\section{Primitive d'une fonction continue}



\begin{definition}Soit $f$ une fonction définie et continue sur un intervalle $I$ et $F$ une fonction dérivable sur cet intervalle. On dit que $F$ est une \textit{primitive} de $f$ sur $I$ si $F'=f$. 

Autrement dit, $F$ est solution de l'équation différentielle $y'=f$.\end{definition}


\begin{example} Pour tout réel $x$, on pose $f(x)=6x^2+4x+3$ et $F(x)=2x^3+2x^2+3x-7$. On a bien pour tout réel $x$, $F'(x)=f(x)$. $F$ est une primitive de $f$ sur $\mathbb{R}$.\end{example}



\begin{theorem}Toute fonction continue sur un intervalle $I$ admet des primitives.\end{theorem}

Ce résultat sera démontré (en partie) dans un prochain chapitre. Patience !

\begin{proposition}Soit $f$ une fonction continue sur un intervalle $I$, $F_1$ et $F_2$ deux primitives de $f$ sur $I$. 

Alors $F_1-F_2$ est constante : deux primitives d'une même fonction continue sur un intervalle différent d'une constante.\end{proposition}

\begin{demonstration} Soit $F_1$ et $F_2$ deux primitives de $f$ sur $I$.
$F_1-F_2$ est dérivable sur $I$ comme différence de fonctions dérivables sur $I$. De plus, pour tout réel $x$ dans $I$
\[ (F_1-F_2)'(x)=F_1'(x)-F_2'(x)=f(x)-f(x)=0.\]
Ainsi, $F_1-F_2$ est de dérivée nulle sur l'intervalle $I$, $F_1-F_2$ est donc constante sur $I$.\end{demonstration}

\begin{proposition}Soit $I$ un intervalle, $x_0\in I$, $y_0 \in \mathbb{R}$ et $f$ une fonction continue sur $I$. Il existe une unique primitive $F$ de $f$ sur $I$ tel que $F(x_0)=y_0$.

L'équation différentielle $y'=f$ ayant pour condition initiale $y(x_0)=y_0$ possède une unique solution.\end{proposition}


\begin{example} Pour tout réel $x$, on pose $f(x)=\dfrac{2x}{x^2+1}$ et $F(x)=\ln(x^2+1)$. 

$F$ est une primitive de $f$ sur $\mathbb{R}$. En effet, pour tout réel $x$, $F'(x)=f(x)$.

On cherche alors l'unique primitive $F_0$ de $f$ qui vaut 5 en 0. Puisque toutes les primitives d'une fonction diffèrent seulement d'une constante, il existe un réel $c$ tel que $F_0=F+c$ et donc, pour tout réel $x$, \[F_0(x)=\ln(x^2+1)+c.\] 

Puisque $F_0(0)=5$, on a alors $\ln(0^2+1)+c=5$, soit $c=5$. Pour tout réel $x$, on a donc \[F_0(x)=\ln(x^2+1)+5.\]\vspace{-0.5cm}\end{example}



\begin{proposition}Primitives usuelles
\vskip10pt
\renewcommand{\arraystretch}{2}
\begin{tabularx}{\linewidth}{|X|X|X|}

\hline
Fonction $f$ & \textbf{UNE} Primitive $F$ & Intervalle $I$ \\
\hline
$x \mapsto x^n$, $n\in \mathbb{N}$ & $x\mapsto \dfrac{x^{n+1}}{n+1}$ & $\mathbb{R}$ \\
\hline
$x \mapsto \dfrac{1}{x^n}$, $n\in \mathbb{N}\setminus\{0;1\}$ & $x\mapsto -\dfrac{1}{(n-1)x^{n-1}}$ & $]-\infty;0[$ et $]0;+\infty[$ \\
\hline
$x\mapsto \dfrac{1}{\sqrt{x}}$ & $2\sqrt{x}$ & $]0;+\infty[$ \\
\hline
$x \mapsto \e^{ax+b}$, $a\in \mathbb{R}^*$, $b\in \mathbb{R}$ & $x\mapsto \dfrac{\e^{ax+b}}{a}$ & $\mathbb{R}$ \\
\hline
$x \mapsto \dfrac{1}{x}$ & $x\mapsto \ln(x)$ & $]0;+\infty[$ \\
\hline
\end{tabularx}\end{proposition}

Le conseil de la leçon : Une fois que vous avez détermine une primitive, dérivez-la ! Vous devez obtenir la fonction de départ. Il est bien plus facile de dériver que de primitiver.

\newpage
\begin{proposition}Soit $f$ et $g$ deux fonctions continues sur un intervalle $I$, $\lambda$, $F$ une primitive de $f$ et $G$ une primitive de $g$. Alors,
\begin{itemize}
\item $F+G$ est une primitive de $f+g$ ;
\item $\lambda\,F$ est une primitive de $\lambda\,f$.
\end{itemize}\end{proposition}

\begin{example}Une primitive de la fonction $f:x\mapsto \e^{2x-1}+5x^2+\dfrac{3}{x^2}$ sur $I=]0;+\infty[$ est la fonction
\[F:x\mapsto \dfrac{\e^{2x-1}}{2}+\dfrac{5}{3}x^3-\dfrac{3}{x}.\]\end{example}



\begin{proposition}L'essentiel : Soit $f$ une fonction dérivable sur un intervalle $I$ et $g$ une fonction dérivable sur un intervalle $J$ telle que pour tout réel $x\in J$, $g(x)\in I$.

$f \circ g$ est une primitive de $g' \times (f' \circ g)$ sur $J$.\end{proposition}

Certaines "formes" de fonction sont à reconnaître pour en calculer les primitives. 

\begin{itemize}
\item Une primitive d'une fonction de la forme $-\dfrac{u'}{u^2}$ où $u$ est une fonction qui ne s'annule pas est $\dfrac{1}{u}$.
\vskip10pt
\item Une primitive d'une fonction de la forme $\dfrac{u'}{2\sqrt{u}}$ où $u$ est une fonction strictement positive est $\sqrt{u}$.
\vskip10pt
\item Une primitive d'une fonction de la forme $u'\e^u$ est $\e^u$.
\vskip10pt
\item Une primitive d'une fonction de la forme $\dfrac{u'}{u}$ où $u$ est une fonction strictement positive est $\ln(u)$.
\end{itemize}


\begin{example} Pour tout réel $x$, on pose $f(x)=(2x+1)\e^{x^2+x-2}$.

Posons, pour tout réel $x$, $u(x)=x^2+x-2$. On remarque alors que $f(x)=u'(x) \times \e^{u(x)}$.

Une primitive de $f$ est donc la fonction $F$ définie pour tout réel $x$ par $F(x)=\e^{u(x)}=\e^{x^2+x-2}$.\end{example}

\begin{example}Pour tout réel $x$, on pose $f(x)=\dfrac{2\e^{2x}}{1+\e^{2x}}$. Posons, pour tout réel $x$, $u(x)=1+\e^{2x}$. 

Pour tout réel $x$, $u(x)>0$. Par ailleurs, on remarque alors que pour tout réel $x$, $f(x)=\dfrac{u'(x)}{u(x)}$. 

Une primitive de $f$ est donc la fonction $F$ définie pour tout réel $x$, par $F(x)=\ln(u(x))=\ln(1+\e^{2x})$.\end{example}

Malheureusement, certaines fonctions n'admettent pas de primitive pouvant être exprimées à l'aide de fonctions usuelles : il s'agit là d'un théorème démontré par Liouville dans les années 1830.

L'exemple le plus notable est la fonction $x\mapsto \e^{-x^2}$, très utilisée en probabilités et dont les applications dans d'autres domaines se comptent par centaines (citons ainsi la balistique, l'évaluation du quotient intellectuel, le traitement du signal,  le cours de la bourse...).

\section{Equation différentielle  du premier ordre}

\subsection{Equations différentielles homogènes $y'+ay=0$}

\begin{definition}Soit $a$ un réel. L'équation $y'+ay=0$ ayant pour inconnue une fonction dérivable $y$ sur $\mathbb{R}$ s'appelle "équation différentielle \textbf{homogène} du premier ordre, à coefficients constants".\end{definition}

\begin{proposition}Soit $a$ un réel.

Les solutions l'équation $y'+ay=0$ sont les fonctions $f$ définies pour tout réel $x$ par $f(x)=C\e^{-ax}$ où $C$ est un réel quelconque.

De plus, pour tous réels $x_0$ et $y_0$, il existe une unique solution $f_0$ de cette équation différentielle telle que $f(x_0)=y_0$.\end{proposition}

\begin{demonstration} Soit $C$ un réel. Pour tout réel $x$, on pose alors $f(x)=C\e^{-ax}$. $f$ est dérivable sur $\mathbb{R}$ et pour tout réel $x$, $f'(x)=-a\times C\e^{-ax}=-af(x)$. Ainsi, pour tout réel $x$, \[f'(x)+af(x)=-af(x)+af(x)=0.\] $f$ est donc bien solution de l'équation différentielle $y'+ay=0$.

Réciproquement, soit $f$ une solution de l'équation $y'+ay=0$. On a alors, pour tout réel $x$, $f'(x)+a\,f(x)=0$.

Pour tout réel $x$, posons alors $g(x)=f(x)\e^{ax}$. $g$ est dérivable comme produit de fonctions dérivables et, pour tout réel $x$,
\[g'(x)=f'(x)\e^{ax}+a\,f(x) \e^{ax}=\e^{ax}(f'(x)+af(x))=0.\]
Ainsi, $g$ est constante : il existe un réel $C$ telle que, pour tout réel $x$, $g(x)=C$, c'est-à-dire $f(x)\e^{ax}=C$ et donc $f(x)=C\e^{-ax}$.\end{demonstration}

\begin{example} Les solutions de l'équation $y'-2y=0$ sont les fonctions $f:x\mapsto C\e^{2x}$ où $C$ est un réel.

Cherchons l'unique solution $f_0$ de cette équation telle que $f_0(3)=1$. Soit $C$ le réel tel que, pour tout réel $x$, $f_0(x)=C\e^{2x}$.

On a alors $f_0(3)=C\e^{6}$ d'une part et $f_0(3)=1$ d'autre part. Ainsi, $C=\dfrac{1}{\e^6}=\e^{-6}$. 

Finalement, pour tout réel $x$, on a $f_0(x)=\e^{-6}\e^{2x}=\e^{2x-6}$.\end{example}

Ce théorème signifie que si l'on regarde l'ensemble des courbes des fonctions solutions de l'équation $y'=ay$ et que l'on s'intéresse à un point du plan, il existe une unique courbe qui passe par ce point. En particulier, toute fonction solution qui s'annule n'est autre que la fonction nulle.

\begin{center}
\includegraphics[scale=1.2]{equadiff}
\end{center}

\subsection{Equation différentielle $y'+ay=b$, avec $b$ réel}

\begin{definition} Soit $a$ et $b$ deux réels. L'équation $y'+ay=b$ ayant pour inconnue une fonction $y$ dérivable sur $\mathbb{R}$ s'appelle  "équation différentielle du premier ordre, à coefficients constants et à second membre constant".\end{definition}

\begin{proposition}Soit $a$ et $b$ deux réels.

Les solutions de l'équation différentielle $y'+ay=b$ sont les fonctions $f$ définies pour tout réel $x$ par $f(x)=C\e^{-ax}+\dfrac{b}{a}$ où $C$ est un réel quelconque.

De plus, pour tous réels $x_0$ et $y_0$, il existe une unique solution $f_0$ de cette équation telle que $f_0(x_0)=y_0$.\end{proposition}

\begin{demonstration} Cherchons tout d'abord une solution constante $\varphi$ à cette équation. Soit $k$ le réel tel que pour tout réel $x$, $\varphi (x)=k$. $\varphi$ est dérivable sur $\mathbb{R}$ et pour tout réel $x$, $\varphi '(x)=0$.

Puisque $\varphi$ est solution de l'équation $y'+ay=b$, on a alors $ak=b$ c'est-à-dire $k=\dfrac{b}{a}$.

Soit maintenant $f$ une fonction dérivable sur $\mathbb{R}$.

$f$ est solution de l'équation différentielle $y'+ay=b$ si et seulement si $f'+af=b$. Or, $\varphi$ étant également une solution de cette équation, on a donc $b=\varphi'+a\varphi$.

Ainsi, $f$ est solution de l'équation différentielle $y'+ay=b$ si et seulement si $f'+af=\varphi+a\varphi '$, ce qui équivaut à $f'-\varphi'+a(f-\varphi)=0$, ou encore $(f-\varphi)'+a(f-\varphi)=0$.

Ainsi, $f$ est solution de l'équation différentielle $y'+ay=b$ si et seulement si $f-\varphi$ est solution de l'équation différentielle homogène $y'+ay=0$.

Il existe donc un réel $C$ tel que, pour tout réel $x$, $(f-\varphi)(x)=C\e^{-ax}$.

 Ainsi, pour tout réel $x$, $f(x)=C\e^{-ax}+\varphi(x)=C\e^{-ax}+\dfrac{b}{a}$.

\end{demonstration}



\begin{example} On cherche à résoudre l'équation $y'-5y=-2$ .

\begin{itemize}
\item On détermine les solutions de l'équation homogène associée $y'-5y=0$. 

Ce sont les fonctions $x\mapsto C\e^{5x}$ où $C$ est un réel.
\item On cherche une solution constante (égale à $k$) à l'équation de départ $y'-5y=-2$. 

La dérivée d'une fonction constante étant nulle, on a alors $-5k=-2$ et donc $k=\dfrac{2}{5}$.

\item Les fonctions solutions de l'équation $y'=5y-2$ sont donc les fonctions $f$ définies pour tout réel $x$ par $f(x)=C\e^{-5x}+\dfrac{2}{5}$ où $C$ est un réel.
\end{itemize}

On souhaite déterminer l'unique solution $f_0$ de cette équation telle que $f_0(7)=12$. Soit $C$ le réel tel que, pour tout réel $x$, $f_0(x)=C\e^{-5x}+\dfrac{2}{5}$.

On a alors $f_0(7)=C\e^{-35}+\dfrac{2}{5}=12$ et donc $C=\dfrac{58}{5}\e^{+35}$. Ainsi, pour tout réel $x$, $f_0(x)=\dfrac{58}{5}\e^{-5x+35}+\dfrac{2}{5}$.\end{example}


\begin{proposition}[Formule magique] Pour résoudre une équation du type $y'+ay=b$...
\begin{center}
\textbf{SOLUTION GÉNÉRALE = SOLUTION HOMOGÈNE + SOLUTION CONSTANTE}
\end{center}\vspace{-0.5cm}\end{proposition}


\subsection{Equation différentielle $y'+ay=g$, où $g$ est une fonction}

\begin{definition}Soit $a$ un réel non nul et $g$ une fonction définie sur $\mathbb{R}$. L'équation $y'=ay+g$ s'appelle "équation différentielle linéaire du premier ordre avec second membre".\end{definition}

\begin{proposition}Soit $a$ un réel non nul et $g$ une fonction définie sur $\mathbb{R}$. 

Soit $\varphi$ une solution particulière de cette équation. Alors $f$ est solution de l'équation $y'+ay=g$ si et seulement si $f-\varphi$ est solution de l'équation homogène associée $y'+ay=0$.

Autrement dit, toute solution de l'équation $y'+ay=g$ est de la forme $f+\varphi$, où $f$ est solution de l'équation $y'+ay=0$ et $\varphi$ est \textbf{UNE} solution de l'équation $y'+ay=g$.\end{proposition}

\begin{demonstration}La démonstration est en tout point semblable à celle qui a conduit à l'ensemble des fonctions solutions de l'équation $y'+ay=b$. Faisons-la donc à nouveau. Soit $f$ une fonction dérivable sur $\mathbb{R}$.

$f$ est solution de l'équation différentielle $y'+ay=g$ si et seulement si $f'+af=b$. Or, $\varphi$ étant également une solution de cette équation, on a donc $g=\varphi'+a\varphi$.

Ainsi, $f$ est solution de l'équation différentielle $y'+ay=g$ si et seulement si $f'+af=\varphi+a\varphi '$, ce qui équivaut à $f'-\varphi'+a(f-\varphi)=0$, ou encore $(f-\varphi)'+a(f-\varphi)=0$.

Ainsi, $f$ est solution de l'équation différentielle $y'+ay=g$ si et seulement si $f-\varphi$ est solution de l'équation différentielle homogène $y'+ay=0$.

Il existe donc un réel $C$ tel que, pour tout réel $x$, $(f-\varphi)(x)=C\e^{-ax}$ et donc $f(x)=C\e^{-ax}+\varphi(x)$.\end{demonstration}

\begin{example} On considère l'équation différentielle $y'-2y=-6x^2+13$.
\begin{itemize}
\item Les solutions de l'équation homogène associée $y'-2y=0$ sont les fonctions $x\mapsto C\e^{2x}$ où $C\in\mathbb{R}$.
\item La fonction $\varphi:x\mapsto 3x^2+3x-5$ est solution de l'équation différentielle $y'-2y=-6x^2+13$. En effet, $\varphi$ est dérivable et pour tout réel $x$, 
\[\varphi'(x)-2\varphi(x)+6x^2-13=6x+3-2(3x^2+3x-5)+6x^2-13=0.\]
\end{itemize}
Ainsi, l'ensemble des solutions de l'équation différentielle $y'-2y=-6x^2+13$ sont les fonctions $f$ définies pour tout réel $x$ par $f(x)=C\e^{2x}+3x^2+3x-5$, où $C$ est un réel.\end{example}

\begin{proposition}[Formule magique] Pour résoudre une équation du type $y'+ay+g$...
\begin{center}
\textbf{SOLUTION GÉNÉRALE = SOLUTION HOMOGÈNE + SOLUTION PARTICULIÈRE}
\end{center}\vspace{-0.5cm}\end{proposition}


Il existe de nombreux types d'équations différentielles : linéaires, quadratiques, d'ordre divers... Et parmi toutes celles-ci, il en existe peu que l'on sait résoudre.

On compte par exemple les équations de Navier-Stokes qui décrivent les mouvements des fluides newtoniens. Il s'agit d'équation dites "aux dérivées partielles". Les fonctions en jeu utilisent en effet plusieurs variables (de temps et d'espace en l'occurrence) et on dérive alors selon l'une ou l'autre des variables.

Ces équations sont particulières puisque depuis l'an 2000, leur résolution est mise à prix : quiconque permettra une avancée significative dans leur étude recevra, en plus d'un certain prestige auprès de la communauté mathématique, la coquette somme d'un million de dollars. 

6 autres problèmes ont été mis à pris par l'Institut Clay cette même année. Depuis, un seul d'entre eux a été résolu.





\end{document}